\section{实数的\texorpdfstring{$a$}{a}进制展开式}

记号约定:$a\geq 2$.

\begin{definition}{实数的$a$进制展开式}{}
	设$x\in\R$.我们称$x$在底列$\left(a^{n}\right)_{n=0}^{\infty}$下的Cantor级数展开式为$x$的$a$进制展开式,称$a$为底.
\end{definition}

\begin{convention}{实数的$a$进制展开式的小数记号}{}
	当$x\geq 0$时,如果$x$的$a$进制规范展开式为$\lfloor x\rfloor+\sum\limits_{n=0}^{\infty}\dfrac{u_{n+1}}{a^{n+1}}$,则记$p_{0}=\lfloor x\rfloor$,把$p_{0}$的数字写在小数点左侧,再依次把$u_{1},\ldots,u_{n}$写在小数点右侧,所得记号简写为$x=S\ldots$.这里不是说$S\ldots$等于$x$,而是约定$S\ldots$是$x$的一个不足近似,其误差限$\leq\dfrac{1}{a^{n}}$.
	
	当$x<0$时,先对$-x$按上述规则记号表示,再在整个记号前加负号,因此所得记号表示的是$x$的一个$\frac{1}{a^{n}}$-过剩近似.在负对数的传统记法中,通常沿用正数的小数表示方式，并在整数部分上方加横线,横线表示整数部分取负,而小数部分仍按正数处理.例如,$\overline{3}.14=-3+0.14$.
\end{convention}

